\section{Conclusions}
\label{sec:conclusions}

This thesis asked whether a large language model fine-tuned on the largest available single-cell
drug-perturbation atlas learns drug-specific biology. It does not --- and the interesting content is
in the chain of results establishing why, how far the failure can be repaired, and what bounds the
repair.

\paragraph{The evaluation had to be fixed first.} The metric the field standardises on is exploitable:
a zero-information baseline placing every gene at the middle rank scores $\DEdr = 1.000$, above both
the noise ceiling and the model. Under a Dynamic Range Fraction analysis computed across plates, and replicated within plate, only
$\NIR$ --- a discrimination metric, in which the generic response cancels --- has a positive dynamic
range fraction: $+0.635$ across plates and $+0.519$ within plate. All four prediction metrics tested
come out negative --- $-0.163$ to $-0.686$ across plates, $-0.122$ to $-0.500$ within plate ---
inverted, not merely uninformative. These figures are reported descriptively. Both released
calibration artifacts cluster the drug rows as though each row were its own cell line ($1{,}820$ and
$335$ rows, when there are $25$ cell lines), so their bootstrap intervals are far too narrow, and the
one-sided $p$ for $\NIR$ sits at the resampling floor of $1/2001$ in both runs --- the Holm-adjusted
$0.0025$ quoted earlier in this thesis is that floor times the family size of five, not a measured
value. A corrected rerun is pending and no significance claim is made here. This is the most portable result
here, since it applies to any work scored on $\DEdr$-family metrics.

This also settles how the thesis enters the dispute of Section~\ref{sec:lit-dispute}. Miller et
al.~\cite{miller} argue that reports of deep models failing to beat uninformative baselines reflect
miscalibrated evaluation rather than model failure, and on genetic perturbation they support that
with calibrated-metric benchmarks. We port their framework, confirm its methodological claim
emphatically --- metric choice dominates, and only the rank-based metric survives --- and then find
that on drug perturbation the model is at chance \emph{on the metric their own analysis endorses}.
That is the strongest form the sceptical result can take: it cannot be answered with the objection
their paper raises. Their conclusion is not overturned so much as bounded --- it does not extend
from genetic perturbation to this task, and the per-drug-constant finding below suggests why.

\paragraph{The model is drug-blind, and the representation is not the reason.} Four independent
instruments agree that swapping the drug token changes the output by nothing. Yet drug identity is
linearly decodable from the residual stream, rising to $0.88$ by layer 16 with context removed,
while the purified drug subspace's causal variance share stays flat near $3\%$. The model computes an
increasingly explicit drug representation and does not act on it --- a dissociation between what is
encoded and what is used.

How far that can be pushed toward a mechanism is limited, and the limit should be stated plainly.
Two earlier attempts to localise the failure by activation swapping were retracted after their
matched controls were found confounded, and what replaces them
(Section~\ref{sec:res-mechanism}) is a small, replicated, single-layer effect present in one training
arm and absent in another. That supports the modest reading --- the generation readout consults the
drug representation at very low gain --- and does not support a claim about privileged subspaces or
global workspaces, which is offered here as an interpretive frame rather than a result.

\paragraph{The bottleneck appears to be the tokenisation, and it is repairable.} Because a cell
sentence ranks genes by expression, the targets for two different drugs in the same context are
nearly the same string: $82.7\%$ of the top-$200$ genes are shared, so only about a sixth of the
target distinguishes one drug from another. Since the training loss weights every token position
equally, the natural inference is that very little of the gradient can be carrying drug identity ---
but that inference is a \emph{hypothesis} about the optimisation, not a measurement of it. What is
measured here is the overlap of the target gene sets; no token-level loss or gradient attribution was
computed, and the claim should be read at that strength. What follows is that the hypothesis makes
correct predictions, which is the evidence for it. It explains a uniform pattern of failure: an entire $\varepsilon$-ladder of target constructions, from a single
state-matched treated cell to a full pseudobulk collapse, stays in the near-null regime, because all
of them change \emph{what the target is} and none change \emph{what the tokens encode}. (One arm,
the optimal-transport target, does show a weak effect once a sharper stratified test is applied ---
\ci{+0.0263}{+0.0069}{+0.0476} --- but roughly five-fold below what re-encoding achieves, which
is the distinction that matters.) Re-encoding
the target as the drug-specific residual raises the differing fraction to $59\%$ and produces the
first non-null drug effect in this project, with a monotone gradient across swap dissimilarity and a
matched control that stays null.

\paragraph{What that drug use amounts to, and what it does not.} An earlier form of
this paragraph claimed the drug use generalises, on a holdout keyed at the (drug, cell line) level
where pairings never seen in training scored \ci{+0.1002}{+0.0661}{+0.1368} and held-out
\emph{drugs} came out null. Two corrections have since removed that claim. The targets were rebuilt
so the generic is fitted after the split rather than over the holdout, and the comparator was
corrected: the scramble arm the gap was measured against is not a null but a different drug's
learned signature, chosen to point away from the truth. The held-out-drug arm is what exposed it,
by returning a positive result on a split where a positive result is impossible.

Measured against the neutral comparator on the rebuilt targets, trained pairings score $+0.0898$
\nolinebreak[3]$[+0.0382, +0.1415]$, held-out pairings $+0.0332$, and held-out \emph{drugs} $-0.0315$
\nolinebreak[3]$[-0.0836, +0.0206]$. Transfer to held-out combinations is \emph{not} established. Its
interval excludes zero when the bootstrap is clustered on cell line and well
\nolinebreak[3]$[+0.0069, +0.0594]$, and does not when it is clustered on drug and well
\nolinebreak[3]$[-0.0031, +0.0695]$ --- and drug is the axis along which a claim about transferring a
\emph{drug} signature has to generalise. The split is also mostly recombination rather than novelty:
$30$ of its $34$ drugs and $39$ of its $42$ cell lines also appear in the trained arm, and $84\%$ of
its observations sit on those shared drugs. The memorisation premium that earlier carried this
argument goes with it. Pooled it is $+0.0567$ \nolinebreak[3]$[-0.0014, +0.1148]$, cluster-robust
$p = 0.056$; restricted to the $30$ drugs common to both splits it changes sign, to $-0.0606$ with
$p = 0.18$. It is confounded in any case, because the splits are not equally hard: the within-well
split-half precision reference is $0.956$ on the trained split against $0.824$ and $0.836$ on the two
held-out ones, a spread of $0.132$, so gaps measured on them are not directly comparable.

What survives is narrower than what was claimed for it. The comparator-free test --- the slope of
$\NIR$ on the cosine similarity between the real response of the drug named in the prompt and that of
the true one --- is positive on every split: $+0.373$ on trained pairings, $+0.337$ on held-out
pairings, $+0.342$ on held-out drugs. This does \emph{not} show that the model knows held-out
compounds. The slope is fitted over the three scrambled \emph{partner} prompts alone (\texttt{near},
\texttt{orth}, \texttt{opposite}); the prompt naming the true target drug is never a point on the
line, and on the held-out-drug split $82\%$ of the $597$ points name a drug that also appears in the
trained arm. What it shows is that the model responds to the identity of a named trained drug even
when it is scored in a context whose true drug it has never seen --- a real effect, and a much weaker
one than zero-shot drug knowledge. Alongside it, on both held-out splits the model scores at chance
when its output is ranked against other drugs in the same cell line ($\NIR = 0.533$
\nolinebreak[3]$[0.492, 0.575]$ and $0.459$ \nolinebreak[3]$[0.404, 0.514]$), and its output is no
closer to the true drug's response than to other drugs' (own-minus-other cosine $+0.011$ and
$-0.014$, neither interval excluding zero). The map exists and is coarse: what the model has captured
is close to a per-drug average response --- which is also why the lookup below beats it, since the
lookup \emph{is} that average, measured directly.

\paragraph{But a lookup table wins --- and what it beats is not a ceiling.} A per-drug mean residual
measured in other cell lines scores $0.913$; the model reaches $0.537$ overall, and $0.549$ on the
$n = 1192$ conditions on which both are scored. The within-well split-half precision reference is
$0.854$ over all conditions and $0.857$ on that common support, so the lookup does not approach that
reference --- it \emph{exceeds} it, recovering $116\%$ of its headroom above chance on the common
support against the model's $14\%$. The natural reading of those numbers is that the drug's response
is essentially a
per-drug constant, that a conditional model therefore has almost nothing to condition on, and that
the correct formulation of the task is retrieval rather than generation. We initially drew exactly
that conclusion, and it is wrong.

It is wrong because $\NIR$ compresses near the top of its range --- the same class of failure this
thesis opens by documenting, arriving a second time in our own preferred instrument, at the top of
the scale instead of the bottom. Measured in cosine space with measurement noise divided out, the
transfer coefficient is \ci{\Tcoef = 0.553}{0.514}{0.593}, against a simulated null in which the
residual is entirely shared across contexts --- a null this estimator recovers at $1.001$ --- and
with a structure-matched negative control at \ci{-0.008}{-0.036}{+0.020}. So roughly $45\%$
\nolinebreak[3]$[41\%, 49\%]$ of the drug-specific residual is \emph{not shared across contexts}.
That is deliberately weaker than the reading in an earlier draft, which took $1 - \Tcoef$ to be a
drug$\times$cell-line interaction share. It is not one. The estimator pools pairs that differ in
cell line, in dose and in well, and it was run with \texttt{same\_dose\_only} off,
\texttt{loo\_generic} off and a reproducibility threshold of $0.2$; on pairs that vary dose within a
single cell line it returns \ci{\Tcoef = 0.707}{0.632}{0.782}, so nearly $0.3$ of the residual still
fails to transfer with the cell line held fixed. The unshared remainder therefore mixes interaction
with dose, well and estimator slack, and cannot be attributed to a drug$\times$cell-line term
without a design that separates them. What can be said is that the per-drug constant is what the
lookup captures and what the model gropes toward, and that the unshared remainder is reached by
neither.

So the honest statement of what the model does is sharper, and less flattering, than ``it misses the
part that is not shared across contexts''. It fails to recover the \emph{shared} part that a numpy
dictionary obtains for free, and the unshared remainder --- which mixes cell line, dose, well and
estimator, and in which no learnable structure was found --- is reached by no method here, ours
least of all. Two different failures, and collapsing them undersells both.

That conclusion connects this work to an assumption the field has carried since the Connectivity
Map: that a compound induces a characteristic signature stable enough across cellular contexts for
signature-similarity querying to be meaningful. This thesis measures that assumption directly, with
a discrimination metric and a within-well split-half precision reference, at single-cell resolution --- and the modern
finding that a per-perturbation mean is hard to beat may be less a warning about weak baselines than
a restatement of it.

It also places this work as the single-cell, generative instance of DrEval's
critique~\cite{dreval}. DrEval audits drug \emph{sensitivity} prediction and finds that deep models
barely outperform a naive additive predictor of mean drug and cell-line effects. Our
\texttt{drug\_lookup} is that predictor's drug term, and the finding replicates on a
$946$-dimensional transcriptomic response: $0.913$ for the lookup against $0.537$ for the model,
with a within-well split-half precision reference of $0.854$ that the lookup exceeds. Three of DrEval's six pathologies --- pseudoreplication, biased evaluation
metrics, and missing ablation studies --- turn out to structure the entire methodology of this
thesis, which arrived at clustered bootstraps over cell lines, a discrimination metric immune to
between-drug variance, and a token-level scramble ablation independently. That convergence is worth
recording: the controls a critical benchmark demands are the same controls a single negative result
needs in order to mean anything.

\paragraph{Where the opportunity actually is.} Two further measurements bear on where it lies --- one
unresolved, the other provisional pending a corrected rerun --- and between them they redirect what
follows from this work.

The part of the drug-specific signal that is not shared across contexts is large --- roughly $45\%$ ---
and we could not find structure in it. No cell line modulates different drugs in a consistent
direction. Within-line consistency is $-0.0149$ \nolinebreak[3]$[-0.0251, -0.0050]$ against
$+0.0003$ \nolinebreak[3]$[-0.0071, +0.0081]$ across lines, an excess of $-0.0153$; the within-line
interval lies entirely below zero, so there is no positive consistency to find. A direct test of
whether mechanism-matched drugs share a context term with a given cell line returned nothing
detectable.

Those are two negative results and they do not license the conclusion that the unshared part is
irreducible. The second test bounds its own power severely: the mean split-half reliability of
$\kappa$ in this atlas is $0.326$, so the quantity being correlated is barely estimable, and
mechanism is partially confounded with plate, by an amount that was never measured for this test. The
honest statement is that the structure of the unshared part is \emph{unresolved here}, and that
resolving it needs more cells per condition or a plating design that crosses mechanism with plate ---
neither of which is an analysis choice. What can be said without qualification is that a lookup
scoring $0.913$, above the $0.854$ within-well split-half precision reference, already outscores
anything conditional built here, so on seen drugs nothing we built beats retrieval.

But a lookup requires the drug to have been seen. For drugs that have not been, protein-target and
mechanism annotations score above a \emph{count-matched} null
(\ci{+0.1863}{+0.1264}{+0.2462} and \ci{+0.0700}{+0.0198}{+0.1202}), and both survive the
plate-matched re-measurement an earlier draft listed as pending
(\ci{+0.1799}{+0.1138}{+0.2459} and \ci{+0.0663}{+0.0088}{+0.1237}) as well as the stricter
different-plate requirement (\ci{+0.1115}{+0.0448}{+0.1783} and \ci{+0.0604}{+0.0094}{+0.1114});
Table~\ref{tab:fieldattribution} carries all three constructions. Chemical structure is far weaker
and does not survive that requirement (\ci{+0.0183}{-0.0033}{+0.0399}). A different defect is now
the binding one, and it is not yet repaired: the pool the partner drugs were drawn from was not
restricted to training drugs, so a held-out drug could be matched to a partner the model had also
never seen, and the arm therefore does not isolate what is retrievable \emph{from training}. The
code has been corrected and the rerun is pending. Until it lands, these gaps establish that target
and mechanism annotation carry transcriptional signal in this atlas, but they do \emph{not}
establish that unseen-drug generalisation is reachable by target or mechanism retrieval, and that
stronger reading should not be attributed to them.
This at least qualifies a pre-registration made earlier in this project, which inferred that unseen drugs
were unreachable from a gate that tested chemical \emph{structure} --- the channel our own
drug-knowledge specification ranked lowest --- and never read the target annotations sitting in the
atlas's own metadata.

So the thesis does not end in closure. It ends with the opportunity located: on seen drugs a
generative model has nothing to add over retrieval, because the unshared part is unpredictable; on
unseen drugs, where retrieval is impossible by construction, protein-target and mechanism annotation
are the only channels that survive a different-plate null, and --- pending a rerun whose partner pool
is restricted to training drugs --- may be where a model can do something a dictionary cannot. That
is also the setting the benchmark literature calls
Leave-Drugs-Out, and the one that matters for drug design.

\paragraph{A note on method.} The substantive results here rest on controls that were added in
response to specific objections: plate-matched comparison sets after a zero-information baseline
more than doubled its score when the plate was free to vary; a swap-distance sweep after it was
pointed out that scrambling to a
similar drug is a weak test; a removal gate before trusting any ablation null; matched-norm
injection before reading a swap; clustered bootstraps over cell lines rather than conditions; and a
pre-registered interpretation written before the decisive result, whose middle branch is the one
that fired. Several claims were retracted along the way and are struck rather than deleted. In a
literature where the central problem is that uninformative baselines score well, the controls are
not an appendix to the result --- they are most of it.
